Several different simplifications are routinely called ``sparse interaction.''  A model may use
few position pairs, compress the categorical action on each pair, truncate all dependence above
order two, or reduce every edge to one scalar score.  These operations save different resources
and discard different information.  None implies the others.

Graph support answers whether a position pair is represented.  Categorical rank answers how many
independent substitution-contrast directions its edge operator carries.  Interaction order asks
how many sequence coordinates participate irreducibly in one functional component.  A sparse
position graph can attach a full-rank $q\times q$ operator to each surviving edge.  A complete graph
can reuse a one-dimensional message family.  A network built from pairwise messages can generate
higher-order dependence after nonlinear aggregation.

These counterexamples are not pathological.  They identify three distinct approximation budgets:
\begin{equation}
\text{number of edges},\qquad
\text{rank per edge},\qquad
\text{maximum interaction order}.
\label{eq:book-three-complexities}
\end{equation}
A claim of compression should say which budget is controlled and how the discarded component is
measured.

The chapter then confronts the most common compression: replacing an operator $G_{ij}$ by a scalar
such as its norm.  This step is often necessary for contact prediction, visualization, or ordinary
graph spectra.  It is also irreversible.  Sign, singular directions, and the identity of favored
substitutions disappear.  Projection and scalarization do not commute: a large raw norm can come
entirely from one-body structure that vanishes in the canonical pair gauge.

Scalar graphs are therefore not rejected here.  They are located.  A scalar shadow answers a
declared question about the magnitude of an already identified operator.  Trouble begins only when
properties of the shadow are read back as though no categorical information had been discarded.
